% =========================================================
% Finiteness: Compactness — Notes Index
% =========================================================

% ---------------------------------------------------------
% Breadcrumb
% ---------------------------------------------------------
\begin{tcolorbox}[
  colback=gray!6,
  colframe=gray!40,
  arc=2pt,
  left=8pt, right=8pt, top=6pt, bottom=6pt,
  title={\small\textbf{Where You Are in the Journey}},
  fonttitle=\small\bfseries
]
\begin{center}
\small
$\mathbb{N}$, $\mathbb{Z}$, $\mathbb{Q}$, $\mathbb{R}$
$\;\to\;$ Real Analysis
$\;\to\;$ Distance (Metric Spaces)
$\;\to\;$ Nearness (Topology)
$\;\to\;$ \textbf{Finiteness (Compactness)}
$\;\to\;$ Control (Continuity)
$\;\to\;$ $\cdots$
\end{center}

\medskip
\noindent\textbf{How we got here.}
Topology gave us the open-set language for nearness and convergence.
But open covers can be infinite, and infinite witnesses are difficult
to control.
Compactness is the structural property that collapses infinite open
covers to finite ones --- the first time the project acquires
genuinely global power over analytic behavior.

\medskip
\noindent\textbf{What this chapter builds.}
The Heine--Borel theorem ($K \subseteq \mathbb{R}^n$ compact
$\Leftrightarrow$ closed and bounded), the equivalence of
open-cover and sequential compactness, the Lebesgue covering lemma,
nested compact sets, and the first major consequences:
continuous images of compact sets are compact,
continuous functions on compact sets attain their extreme values,
and every continuous function on a compact set is uniformly continuous.

\medskip
\noindent\textbf{Where this leads.}
Uniform continuity on compact domains is precisely the hypothesis
that makes Riemann sums converge.
Compactness is also the key hypothesis in the Arzel\`{a}--Ascoli theorem,
which characterizes compact sets in function spaces and underpins
much of functional analysis.
\end{tcolorbox}
\vspace{1em}

% ---------------------------------------------------------
% Structural Role
% ---------------------------------------------------------
\begin{tcolorbox}[
  colback=gray!6,
  colframe=gray!40,
  arc=2pt,
  left=8pt, right=8pt, top=6pt, bottom=6pt,
  title={\small\textbf{Structural Role: Finiteness — Infinite Behavior with Finite Control}},
  fonttitle=\small\bfseries
]
Compactness is the global principle of analysis.
It asserts that any infinite open covering of a set
admits a finite subcovering.

In~$\mathbb{R}^n$, this coincides with being closed and bounded
--- the Heine--Borel theorem.
In metric spaces, sequential compactness gives an equivalent
formulation: every sequence in a compact set has a subsequence
converging inside the set.

The power of compactness is that it transforms infinite complexity
into finite manageability.
It guarantees:

\begin{itemize}
  \item Limits exist where boundedness alone is insufficient,
  \item Continuous functions attain their extreme values,
  \item Infinite behavior is globally constrained by finitely many
        witnesses.
\end{itemize}

Compactness is the finiteness principle of analysis.
It is where the completeness of~$\mathbb{R}$, the topology of open
sets, and the sequential behavior of earlier chapters converge
into a single structural theorem.
\end{tcolorbox}
\vspace{1em}

% ---------------------------------------------------------
% Structural Roadmap
% ---------------------------------------------------------
\subsection*{Structural Roadmap}

\begin{center}
\textbf{Definitions $\longrightarrow$ Main Theorems
$\longrightarrow$ Consequences and Logical Structure}
\end{center}

The development follows the definition--theorem--structure
architecture used throughout, now applying the open-set framework
to establish global finiteness results.
The global progression is:

\begin{enumerate}
  \item Open-cover compactness: definition and first properties
  \item Sequential compactness and the Bolzano--Weierstrass connection
  \item Equivalence of open-cover and sequential compactness
  \item Heine--Borel theorem: compact $\Leftrightarrow$ closed and
        bounded in~$\mathbb{R}^n$
  \item Lebesgue covering lemma
  \item Nested compact sets and intersection properties
  \item Compactness under continuous maps (first consequences
        for continuity chapter)
\end{enumerate}

\begin{remark*}[Primary sources]
Abbott, \textit{Understanding Analysis}, Chapter~3;
Pugh, \textit{Real Mathematical Analysis}, Chapter~2 §5--6;
Rudin, \textit{Principles of Mathematical Analysis}, Chapter~2.
\end{remark*}

% ---------------------------------------------------------
% Status
% ---------------------------------------------------------
\vspace{1em}
\begin{tcolorbox}[
  colback=gray!6, colframe=gray!40, arc=2pt,
  left=8pt, right=8pt, top=6pt, bottom=6pt,
  title={\small\textbf{Status: Planned}},
  fonttitle=\small\bfseries
]
\begin{center}\Large\bfseries Coming Soon\end{center}
\vspace{6pt}
\noindent Notes, proofs, and exercises will appear here in a future revision.
\end{tcolorbox}

% Future sections (uncomment as content is written):
% \input{volume-iii/analysis/compactness/notes/notes-open-cover}
% \input{volume-iii/analysis/compactness/notes/notes-sequential}
% \input{volume-iii/analysis/compactness/notes/notes-heine-borel}
% \input{volume-iii/analysis/compactness/notes/notes-lebesgue-lemma}
% \input{volume-iii/analysis/compactness/notes/notes-consequences}
